\documentclass[10pt]{article}
\usepackage[margin=0.75in,top=0.6in,bottom=0.6in]{geometry}
\usepackage{amsmath,amssymb,booktabs,microtype,xcolor,url}
\usepackage[colorlinks=true,linkcolor=blue!60!black,urlcolor=blue!60!black,pdftitle={FRE-GY-9743 Assignment 4},pdfauthor={Raj Pawar}]{hyperref}
\newcommand{\code}[1]{\texttt{#1}}
\newcommand{\E}{\mathbb{E}}
\setlength{\parindent}{0pt}\setlength{\parskip}{4pt}\emergencystretch=3em\pagestyle{empty}
\makeatletter\renewcommand{\section}{\@startsection{section}{1}{\z@}{8pt plus 2pt}{3pt}{\normalfont\large\bfseries}}
\renewcommand{\subsection}{\@startsection{subsection}{2}{\z@}{5pt plus 1pt}{1pt}{\normalfont\normalsize\bfseries}}\makeatother
\begin{document}
{\Large\bfseries FRE-GY-9743 Assignment 4: Futures Settlement and IFR Yield Curves}\\[1pt]
Raj Pawar, NYU Tandon School of Engineering, October 18, 2026.\\
Code: \href{https://github.com/pawarraj8888/FRE-GY-9743-Assignments-4/tree/Raj_Pawar_Branch}{github.com/pawarraj8888/FRE-GY-9743-Assignments-4}, branch \code{Raj\_Pawar\_Branch}.
Project page: \href{https://pawarraj8888.github.io/FRE-GY-9743-Assignments-4/}{pawarraj8888.github.io/FRE-GY-9743-Assignments-4}.
\vspace{2pt}\hrule

\section*{Part I: why the OTC price carries a discount factor and the futures price does not}
Both trades move cash every day, and in both cases the amount is the same: today's mark-to-market minus what has already been
posted. The difference is what that cash \emph{is}.

\textbf{Under a CSA the daily transfer is collateral, not settlement.} The swap or forward is still a claim on its payoff $V(T)$ at $T$.
What I post or receive each day is the other side's money parked against an unrealised mark: whoever holds it pays interest on it at
the collateral rate $r_C$, and it flows back when the mark reverses or the trade matures. Nothing about the payoff has been paid early,
so the value is still the present value of $V(T)$, and under perfect collateralisation the notes' equation (16) gives
\[
V(t) = \E^{Q}_t\!\Bigl[e^{-\int_t^T r_C(u)\,du}\,V(T)\Bigr].
\]
The intuition I find clearest: being in the money by $V(t)$ and holding $V(t)$ of the other side's collateral is the same as having been
lent $V(t)$ today against $V(T)$ at maturity, and a loan from $t$ to $T$ at $r_C$ is exactly the factor $e^{-\int_t^T r_C}$. The daily
collateral calls do not change \emph{when} the payoff is realised; they only make sure the money is there if someone defaults.

\textbf{On the exchange the daily transfer is a settlement.} At the close the clearing house pays me $F(t_{k+1},T)-F(t_k,T)$ outright, pays no
interest on it, and re-strikes my contract at the new price so that it is worth exactly zero again. Nothing is ever owed for longer than one
margining period. Writing the zero-value condition for that single period, with the overnight rate for the day fixed at its start so that
$D(t_k,t_{k+1})$ is known at $t_k$,
\[
0 = \E^{Q}_{t_k}\!\bigl[D(t_k,t_{k+1})\,(F(t_{k+1},T)-F(t_k,T))\bigr] = D(t_k,t_{k+1})\bigl(\E^{Q}_{t_k}[F(t_{k+1},T)] - F(t_k,T)\bigr),
\]
so the discount factor drops out, $F(t_k,T)=\E^{Q}_{t_k}[F(t_{k+1},T)]$, and chaining the days gives $F(t,T)=\E^{Q}_t[S(T)]$. The
discounting has not disappeared; it is applied one day at a time, over a horizon so short that the rate is known, and each day's bet is
fully paid before the next day's rate is set. A one-day bet that must have zero value is priced the same whatever today's overnight rate
is, because the rate multiplies both sides of the equation. The cumulative variation margin over the life is $F(T,T)-F(t,T)=S(T)-F(t,T)$,
i.e.\ the forward-style payoff received in instalments, which is why no further time value of money appears in the price.

\textbf{Same cash, different ownership.} Variation margin on a future is not a deposit: the day's loss is paid to the winner and belongs
to the winner, so nobody owes interest on it. Collateral under a CSA still belongs to the party that posted it; the party holding it pays
interest at $r_C$ because it is holding someone else's money against a mark that has not been realised. The discount factor shows up exactly
where interest is paid on the posted amount. The cleared-swap market illustrates this: when the clearing houses moved variation margin
from collateralised-to-market to settled-to-market, they kept paying a price alignment amount equal to the old price alignment interest at
$r_C$, precisely so that the swap keeps its collateralised-swap value; only if that payment were dropped would the contract price like a future.

\textbf{What it changes.} With a horizon-long, random discount factor the OTC forward price is a $T$-forward-measure expectation,
$F^{\mathrm{fwd}}(t,T)=\E^{T}_t[S(T)]$ (notes, section 2.3), while the futures price is $\E^{Q}_t[S(T)]$; the gap is the convexity
adjustment, which comes from the covariance between the discount factor and the payoff. If rates were deterministic the discount factor
could be taken out of the OTC expectation too and the two prices would coincide.

\section*{Part II: computational part}
I completed the two \code{TODO} blocks in \path{fixedincomelib/yield_curve/yield_curve_model.py}. The component discount factor is
$D_{\mathrm{local}}(t)=\exp\bigl(-\int_0^t f(s)\,ds\bigr)$, with the integral taken from the component's piecewise-constant interpolator over
curve time measured in ACT/ACT (ISDA) from the value date. The model's discount factor follows the registered reference recursively:
a component without a reference is its own base (reference factor $1$), and a funding component returns
$D_{\mathrm{ref}}(t)\,D_{\mathrm{local}}(t)$, so \code{SOFR-1B-FLAT} is \code{SOFR-1B} times its own spread factor. The interpolator
from Assignment~1 is carried over into the identical skeleton in \path{fixedincomelib/utilities/numerics.py}. In the notebook I created
the 5-year \code{ProductFixedAccruedCashflow} through \code{qfCreateProductFixedAccruedCashflow} with payment on the termination date,
discounted its payment $c=\text{notional}\times\text{accrued}$ on \code{SOFR-1B-FLAT}, and computed the analytic IFR risk with the chain rule
$\partial V/\partial f_i=-V\,w_i$, where $w_i=\partial(\int_0^T f)/\partial f_i$ comes from
\code{gradient\_of\_integrated\_value\_wrt\_ordinate(0,T)} for both components. The notebook's own reference outputs are reproduced exactly.

\begin{center}\small\setlength{\tabcolsep}{6pt}
\begin{tabular}{@{}llrrr@{}}\toprule
Tenor & Date & Curve time & $D_{\mathrm{SOFR\text{-}1B}}$ & $D_{\mathrm{SOFR\text{-}1B\text{-}FLAT}}$\\ \midrule
6M & 2027-04-07 & 0.498630 & 0.985152 & 0.985152 \\
1Y & 2027-10-07 & 1.000000 & 0.970446 & 0.970446 \\
2Y & 2028-10-10 & 2.008840 & 0.940565 & 0.940565 \\
5Y & 2031-10-07 & 5.000000 & 0.852156 & 0.852156 \\
10Y & 2036-10-07 & 10.000644 & 0.708223 & 0.708223 \\
15Y & 2041-10-07 & 15.000000 & 0.582765 & 0.582765 \\
\bottomrule\end{tabular}\end{center}

With the IFR grid $0.030,0.031,\dots,0.039$ on the 1Y to 10Y nodes and a zero funding spread, the two curves agree. The 5-year cashflow pays
on 2031-10-07: $c=507.222222$ (ACT/360), $D=0.852156$, $V=432.232456$. Its
risk to the first five SOFR nodes is $(-432.2325,\ -436.0536,\ -430.7797,\ -429.8641,\ -432.2325)$ per unit of IFR, i.e.\ $-V$ times the length of each node's bucket inside
$[0,T]$ with $T=5.000000$; nodes beyond five years carry zero risk, and the single funding node carries $-TV=-2161.162279$.
\path{tests/test_yield_curve.py} (14 tests) checks the discount factors against a by-hand integral of the step function, flat
extrapolation past the last node, the reference recursion with zero, flat and multi-node funding spreads, the notebook's numbers, the analytic
risk against central bump-and-reval through \code{perturb\_model\_parameter}, serialization and the validation errors;
\path{tests/test_interpolator_pcp.py} (20 tests) covers the interpolator.
\end{document}
